\scanpage{68}
\mspage{8}
% archive foliation: 70
par bouchage des trous correspondants reste
la $=$ (à isom. can. près). Il faudrait,
pour le voir clairement, voir la
relation entre la sphère \inserted{\unsure{bd}} $S(I)$ triponctuée,
associée à la sphère à bord $\Sigma^*$, et ses
paramètres $\ell_i$ $(i\in I)$ comme \unsure{déduite}
\deleted{(\unsure{obtenue}}
de $S(I)$ en enlevant des disques \deleted{que}
$\unsure{D}(\ell_i)$ \inserted{de centre $s_i$ ?}\deleted{qui} dont chacun ne dépend que
des $\ell_i$ correspondants, pas des autres
$\ell_j$.
\margin{?}

Ce point étant supposé acquis, ou quelque
variante, il semble donc bien que
\inserted{(du moins dans le cas hyperbolique orienté)}
pour une surface conforme ponctuée
donnée, \inserted{de type $(g,\nu)$} les familles \inserted{\unsure{admiss.} $(D_i)$} des disques autour
des $a_i$ forment un espace homéom.
à $\mathbb R^\nu$ (où $\nu=\operatorname{card} I$), i.e. à celui
des systèmes d'ombres possibles ! Il
est donc très tentant de faire la

\uline{Conjecture} Pour tt \unsure{surface} conforme
ponctuée (ou \deleted{tout} au moins dans le cas
hyperbolique), l'application \deleted{que}
\[
(D_i)\longmapsto(o(D_i))
\]
